%%%%%%%%%%%%%%%%%%%%%%%%%%%%%%%%%%%%%%%%
%
% $ Autor: Gruppe 07 $
% $ Projekt: Magic Faucet Fountain $
% $ Pfad: LaTeXTemplate/Montageanleitung/Chapters/de/Montageablauf.tex $
% $ Kurzbeschreibung: Dieses Dokument zeigt den Montageablauf des Magic Faucet Fountain. $
%
% !TeX spellcheck = de_DE/GB
% !TeX program = pdflatex
% !BIB program = biber/bibtex
% !TeX encoding = utf8
%
%%%%%%%%%%%%%%%%%%%%%%%%%%%%%%%%%%%%%%%%

\chapter{Montageablauf}

Bevor Sie mit der Montage beginnen, entnehmen Sie bitte alle vormontierten Bauteile aus der Verpackung und prüfen Sie diese auf Vollständigkeit und Beschädigungen. Wählen Sie einen sauberen, ebenen und stabilen Arbeitsplatz. Eine weiche Unterlage (z. B. eine Decke) schützt die Oberflächen der Bauteile während der Montage.

\begin{enumerate}
	\item Stellen Sie den Blumentopf auf die vorbereitete Fläche.
	\begin{figure}[H]
		\centering
		\resizebox{0.75\textwidth}{!}{%
			\begin{circuitikz}
				\tikzstyle{every node}=[font=\LARGE]
				\draw [line width=1pt, short] (-13.75,18) -- (-11.25,4.25);
				\draw [line width=1pt, short] (-11.25,4.25) -- (1.25,4.25);
				\draw [line width=1pt, short] (1.25,4.25) -- (3.75,18);
				\draw [line width=1pt, short] (-13.75,18) -- (-14,18);
				\draw [line width=1pt, short] (3.75,18) -- (4,18);
				\draw [line width=1pt, short] (-14,18) -- (-14.25,19.25);
				\draw [line width=1pt, short] (-14.25,19.25) -- (4.25,19.25);
				\draw [line width=1pt, short] (4,18) -- (4.25,19.25);
				\draw [line width=1pt, short] (-13.75,18) -- (3.75,18);
				\draw [line width=1pt, short] (-27.5,8) -- (-12,8);
				\draw [line width=1pt, short] (2,8) -- (18.75,8);
				\draw [line width=1pt, short] (-20,8) -- (-25,1.75);
				\draw [line width=1pt, short] (11.25,8) -- (6.25,1.75);
				\draw [ line width=1pt ] (-27.5,28) rectangle (18.75,1.75);
			\end{circuitikz}
		}%
		\caption{Blumentopf positionieren}
		\label{Blumentopfplatzieren}
	\end{figure}
	\item Nehmen Sie den vormontierten Deckel zur Hand. Die Pumpe, der Sensor und die Kabel sind bereits im Deckel montiert und verbunden.
	\begin{figure}[H]
		\centering
		\resizebox{0.75\textwidth}{!}{%
\begin{circuitikz}
	\tikzstyle{every node}=[font=\LARGE]
	\draw [ line width=1pt ] (-10,6.75) rectangle (-7.5,5.5);
	\draw [ line width=1pt ] (-10,6) rectangle (-11,5.5);
	\draw [line width=1pt, short] (-15,18) -- (-1.25,18);
	\draw [line width=1pt, short] (-15,18) -- (-15.25,18.5);
	\draw [line width=1pt, short] (-15.25,18.5) -- (-16,18.5);
	\draw [line width=1pt, short] (-16,18.5) -- (-16,18.75);
	\draw [line width=1pt, short] (-16,18.75) -- (-0.25,18.75);
	\draw [line width=1pt, short] (-0.25,18.75) -- (-0.25,18.5);
	\draw [line width=1pt, short] (-0.25,18.5) -- (-1,18.5);
	\draw [line width=1pt, short] (-1.25,18) -- (-1,18.5);
	\draw [line width=1pt, short] (-8.5,6.75) -- (-8.5,18);
	\draw [line width=1pt, short] (-7.75,18) -- (-7.75,6.75);
	\draw [line width=1pt, short] (-8.5,18.75) -- (-8.5,24);
	\draw [line width=1pt, short] (-7.75,18.75) -- (-7.75,24);
	\draw [line width=1pt, short] (-8.75,24) -- (-7.5,24);
	\draw [line width=1pt, short] (-7.5,24) -- (-7.5,24.5);
	\draw [line width=1pt, short] (-7.5,24.5) -- (-8.75,24.5);
	\draw [line width=1pt, short] (-8.75,24.5) -- (-8.75,24);
	\draw [line width=1pt, short] (-7.75,24.5) .. controls (-7.75,26) and (-7.25,26.25) .. (-9.5,26.5);
	\draw [line width=1pt, short] (-8.5,24.5) .. controls (-8.5,25.5) and (-8.25,25.75) .. (-9.5,25.75);
	\draw [line width=1pt, short] (-9.5,25.75) .. controls (-9.75,25.5) and (-9.75,25) .. (-10.5,25);
	\draw [line width=1pt, short] (-10.5,25) .. controls (-11.25,25) and (-12,25) .. (-12,26);
	\draw [line width=1pt, short] (-12,26) -- (-12.5,26);
	\draw [line width=1pt, short] (-12.5,26.5) .. controls (-12,26.5) and (-12,26.5) .. (-11.75,26.75);
	\draw [line width=1pt, short] (-11.75,26.75) -- (-10.25,26.75);
	\draw [line width=1pt, short] (-10.25,26.75) .. controls (-9.5,26.5) and (-9.75,26.5) .. (-9.5,26.5);
	\draw [line width=1pt, short] (-11.5,26.75) -- (-11.5,27);
	\draw [line width=1pt, short] (-11.5,27) -- (-10.5,27);
	\draw [line width=1pt, short] (-10.5,27) -- (-10.5,26.75);
	\draw [ line width=1pt ] (-11.25,27) rectangle (-10.75,27.5);
	\draw [ line width=1pt ] (-11.75,27.75) rectangle (-10.25,27.5);
	\draw [ line width=1pt ] (-12.75,26.75) rectangle (-12.5,25.75);
	\draw [line width=1pt, short] (-5,18) -- (-5,6);
	\draw [line width=1pt, short] (-5,6) -- (-2.5,6);
	\draw [line width=1pt, short] (-2.5,6) -- (-2.5,18);
	\draw [line width=1pt, short] (-5,6.75) -- (-2.5,6.75);
	\draw [line width=1pt, short] (-7.5,6.25) .. controls (-4.5,5.25) and (-0.75,3) .. (0.75,4.5);
	\draw [line width=1pt, short] (-2.5,16) .. controls (-1.75,15.75) and (-1.75,15.75) .. (-1.25,16);
	\draw [line width=1pt, short] (-1.25,16) .. controls (2.25,17.5) and (1.5,19.5) .. (4,19.75);
	\draw [ line width=1pt ] (-27.5,3) rectangle (18.75,29.25);
	\draw [line width=1pt, short] (0.75,4.5) .. controls (4,7.5) and (-0.25,10.25) .. (-1.25,16);
	\draw [ line width=1pt ] (8.75,20.5) rectangle (13.75,19);
	\draw [line width=1pt, short] (4.5,19.75) -- (8.75,19.75);
	\draw [line width=1pt, short] (4,21) .. controls (3.5,19.5) and (4.5,19.5) .. (4,18.5);
	\draw [line width=1pt, short] (4.5,21) .. controls (4,20) and (5,19.25) .. (4.5,18.5);
\end{circuitikz}
}%
		\caption{Deckel entnehmen}
		\label{Deckel}
	\end{figure}
	\item Legen Sie die Kabel vorsichtig in die dafür vorgesehene seitliche Einbuchtung am oberen Rand des Blumentopfs. Achten Sie darauf, dass die Kabel nicht eingeklemmt oder übermäßig geknickt werden.
	\begin{figure}[H]
		\centering
		\resizebox{0.75\textwidth}{!}{%
		\begin{circuitikz}
			\tikzstyle{every node}=[font=\LARGE]
			\draw [line width=1.2pt, short] (-41.5,3.25) -- (-41.25,0.75);
			\draw [line width=1.2pt, short] (-40.25,6.75) -- (-41,6.5);
			\draw [line width=1.2pt, short] (-39.5,6) -- (-39.5,4.25);
			\draw [line width=1.2pt, short] (-39.5,6) -- (-38.25,6.25);
			\draw [line width=1.2pt, short] (-40.25,4.5) .. controls (-40,4.25) and (-39.75,4) .. (-39.5,4.5);
			\draw [line width=1.2pt, short] (-41,4.25) -- (-40.25,4.5);
			\draw [line width=1.2pt, short] (-41,4.25) .. controls (-40.75,3) and (-39.5,3.5) .. (-39.5,4.25);
			\draw [line width=1.2pt, short] (-39.5,5.75) .. controls (-43,6) and (-45,5) .. (-46.75,0.75);
			\draw [line width=1.2pt, short] (-39.5,5) .. controls (-42.75,5.25) and (-44,4.75) .. (-46,0.75);
			\draw [line width=1.2pt, short] (-42.5,7.5) -- (-42.5,5.5);
			\draw [line width=1.2pt, short] (-42.5,4.25) -- (-42.5,4.75);
			\draw [line width=1.2pt, short] (-41,6.5) -- (-41,5.75);
			\draw [line width=1.2pt, short] (-41,5) -- (-41,4.25);
			\draw [line width=1.2pt, short] (-42.5,4.25) .. controls (-42.25,2.75) and (-38,2) .. (-35,1.75);
			\draw [line width=1.2pt, short] (-42.5,7.5) .. controls (-42.5,7) and (-41.5,6.5) .. (-41,6.5);
			\draw [line width=1.2pt, short] (-39.5,6) .. controls (-38,5.5) and (-38,5.5) .. (-35,5.25);
			\draw [line width=1.2pt, short] (-42.5,7.5) .. controls (-42.5,8.25) and (-40.75,8.75) .. (-39.25,9);
			\draw [line width=1.2pt, short] (-40.25,6.75) .. controls (-40.75,6.75) and (-41.5,7) .. (-41.5,7.5);
			\draw [line width=1.2pt, short] (-41.5,7.5) .. controls (-41.5,8.25) and (-39.5,8.5) .. (-37.25,9);
			\draw [line width=1.2pt, short] (-38.25,6.25) .. controls (-37,5.75) and (-35.75,5.75) .. (-35,5.75);
			\draw [line width=1.2pt, short] (-40.25,6.75) -- (-40.25,5.75);
			\draw [line width=1.2pt, short] (-40.25,5) -- (-40.25,4.5);
			\draw [ line width=1.2pt ] (-48,9) rectangle (-35,0.75);
		\end{circuitikz}
		}%
		\caption{Kabel-Einbuchtung}
		\label{fig:kabeleinbuchtung}
	\end{figure}
	\item Setzen Sie nun den Deckel passgenau auf den Blumentopf. Eine zusätzliche Befestigung ist nicht notwendig – der Deckel liegt stabil auf.
	\begin{figure}[H]
		\centering
		\resizebox{0.75\textwidth}{!}{%
			\begin{circuitikz}
				\tikzstyle{every node}=[font=\LARGE]
				\draw [line width=1pt, short] (-15,1.75) -- (-17.5,11.75);
				\draw [line width=1pt, short] (-17.5,11.75) -- (-17.75,11.75);
				\draw [line width=1pt, short] (-17.75,11.75) -- (-18.5,14.25);
				\draw [line width=1pt, short] (-18.5,14.25) -- (9.75,14.25);
				\draw [line width=1pt, short] (9.75,14.25) -- (9,11.75);
				\draw [line width=1pt, short] (9,11.75) -- (8.75,11.75);
				\draw [line width=1pt, short] (8.75,11.75) -- (6.25,1.75);
				\draw [line width=1pt, short] (-17.5,11.75) -- (8.75,11.75);
				\draw [line width=1pt, short] (-5,14.25) -- (-5,23);
				\draw [line width=1pt, short] (-3.75,23) -- (-3.75,14.25);
				\draw [line width=1pt, short] (-16.75,23) -- (8,23);
				\draw [line width=1pt, short] (8,23) -- (8.25,23.5);
				\draw [line width=1pt, short] (-16.75,23) -- (-17,23.5);
				\draw [line width=1pt, short] (-17,23.5) -- (-17.5,23.5);
				\draw [line width=1pt, short] (-17.5,23.5) -- (-17.5,23.75);
				\draw [line width=1pt, short] (-17.5,23.75) -- (8.75,23.75);
				\draw [line width=1pt, short] (8.25,23.5) -- (8.75,23.5);
				\draw [line width=1pt, short] (8.75,23.75) -- (8.75,23.5);
				\draw [line width=1pt, short] (0,23) -- (0,14.25);
				\draw [line width=1pt, short] (3.75,23) -- (3.75,14.25);
				\draw [line width=1pt, short] (-5,23.75) -- (-5,28);
				\draw [line width=1pt, short] (-3.75,23.75) -- (-3.75,28);
				\draw [ line width=1pt ] (-26.25,1.75) rectangle (17.5,28);
				\draw [line width=1pt, short] (-11.25,15.5) -- (-13.75,16.75);
				\draw [line width=1pt, short] (-11.25,15.5) -- (-8.75,16.75);
				\draw [line width=1pt, short] (-8.75,16.75) -- (-10,16.75);
				\draw [line width=1pt, short] (-13.75,16.75) -- (-12.5,16.75);
				\draw [line width=1pt, short] (-12.5,16.75) -- (-12.5,21.75);
				\draw [line width=1pt, short] (-12.5,21.75) -- (-10,21.75);
				\draw [line width=1pt, short] (-10,21.75) -- (-10,16.75);
			\end{circuitikz}
		}%
		\caption{Deckel Positionierung}
		\label{DeckelimTopf}
	\end{figure}
	\item Stellen Sie sicher, dass der E-Kasten mit den bereits verbundenen Kabeln frei positioniert und nicht eingeklemmt ist.
	\begin{figure}[H]
		\centering
		\resizebox{0.75\textwidth}{!}{%
			\begin{circuitikz}
				\tikzstyle{every node}=[font=\LARGE]
				\draw [ line width=1pt ] (-22.5,26.75) rectangle (-6.25,5.5);
				\draw [ line width=1pt ] (13.75,5.5) rectangle (-2.5,26.75);
				\draw [line width=1pt, short] (-2.5,21.75) .. controls (8,22.75) and (9.5,22.5) .. (13.75,10.75);
				\draw [line width=1pt, short] (-2.5,21) .. controls (9.75,22.75) and (8.5,20) .. (13.75,10);
				\draw [line width=1pt, short] (-22.5,21.75) .. controls (-9,29) and (-20.5,28) .. (-6.25,10.75);
				\draw [line width=1pt, short] (-22.5,21) .. controls (-10,29) and (-19.75,25) .. (-6.25,10);
				\draw [line width=1pt, short] (-20,8) -- (-17.5,10.5);
				\draw [line width=1pt, short] (-20,10.5) -- (-17.5,8);
				\draw [line width=1pt, short] (0,9.25) -- (1.25,8);
				\draw [line width=1pt, short] (1.25,8) -- (2.5,10.5);
				\draw [ line width=1pt ] (-27.5,29.25) rectangle (18.75,3);
			\end{circuitikz}
		}%
		\caption{Kabellage}
		\label{Kabel}
	\end{figure}
	\item Füllen Sie ca. 5 Liter Wasser in den Blumentopf, bis die Pumpe vollständig bedeckt ist.
	\begin{figure}[H]
		\centering
		\resizebox{0.75\textwidth}{!}{%
			\begin{circuitikz}
				\tikzstyle{every node}=[font=\LARGE]
				\draw [line width=1pt, short] (-15,1.75) -- (-17.5,11.75);
				\draw [line width=1pt, short] (-17.5,11.75) -- (-17.75,11.75);
				\draw [line width=1pt, short] (-17.75,11.75) -- (-18.5,14.25);
				\draw [line width=1pt, short] (-18.5,14.25) -- (9.75,14.25);
				\draw [line width=1pt, short] (9.75,14.25) -- (9,11.75);
				\draw [line width=1pt, short] (9,11.75) -- (8.75,11.75);
				\draw [line width=1pt, short] (8.75,11.75) -- (6.25,1.75);
				\draw [line width=1pt, short] (-17.5,11.75) -- (8.75,11.75);
				\draw [line width=1pt, short] (-5,14.25) -- (-5,28);
				\draw [line width=1pt, short] (-3.75,28) -- (-3.75,14.25);
				\draw [line width=1pt, short] (-27.5,17.75) -- (-17.25,20.75);
				\draw [line width=1pt, short] (-17.25,20.75) -- (-21.25,28);
				\draw [line width=1pt, short] (-18.5,23) .. controls (-14.75,20.75) and (-14.5,20) .. (-12.75,14.25);
				\draw [line width=1pt, short] (-19.5,24.75) .. controls (-14,22.75) and (-12.75,22) .. (-9.75,14.25);
				\draw [line width=1pt, short] (-20.75,27) .. controls (-13.5,25.5) and (-10.75,24.25) .. (-6.75,14.25);
				\draw [ line width=1pt ] (-27.5,28) rectangle (18.75,1.75);
			\end{circuitikz}
		}%
		\caption{Wasserbefüllung}
		\label{Wasserbefüllen}
	\end{figure}
	\item Platzieren Sie die Dekosteine gleichmäßig auf dem Deckel. Diese dienen ausschließlich dekorativen Zwecken.
	\begin{figure}[H]
		\centering
		\resizebox{0.75\textwidth}{!}{%
			\begin{circuitikz}
				\tikzstyle{every node}=[font=\LARGE]
				\draw [line width=1pt, short] (-15,1.75) -- (-17.5,11.75);
				\draw [line width=1pt, short] (-17.5,11.75) -- (-17.75,11.75);
				\draw [line width=1pt, short] (-17.75,11.75) -- (-18.5,14.25);
				\draw [line width=1pt, short] (-18.5,14.25) -- (9.75,14.25);
				\draw [line width=1pt, short] (9.75,14.25) -- (9,11.75);
				\draw [line width=1pt, short] (9,11.75) -- (8.75,11.75);
				\draw [line width=1pt, short] (8.75,11.75) -- (6.25,1.75);
				\draw [line width=1pt, short] (-17.5,11.75) -- (8.75,11.75);
				\draw [line width=1pt, short] (-5,14.25) -- (-5,28);
				\draw [line width=1pt, short] (-3.75,28) -- (-3.75,14.25);
				\draw [line width=1pt, short] (-27.5,17.75) -- (-17.25,20.75);
				\draw [line width=1pt, short] (-17.25,20.75) -- (-21.25,28);
				\draw [ line width=1pt ] (-15,20.75) ellipse (0.75cm and 1.25cm);
				\draw [ line width=1pt ] (-11.5,16.75) ellipse (1.25cm and 0.75cm);
				\draw [ line width=1pt ] (-12.5,20.25) ellipse (0.5cm and 0.75cm);
				\draw [ line width=1pt ] (-14,24.5) ellipse (0.75cm and 0.5cm);
				\draw [ line width=1pt ] (-12.25,23) ellipse (1.25cm and 0.5cm);
				\draw [ line width=1pt ] (-18,26) ellipse (0.5cm and 1cm);
				\draw [ line width=1pt ] (-16.5,24.25) ellipse (1cm and 0.5cm);
				\draw [ line width=1pt ] (-16.25,23) ellipse (0.25cm and 0.25cm);
				\draw [ line width=1pt ] (-27.5,28) rectangle (18.75,1.75);
			\end{circuitikz}
		}%
		\caption{Befüllen mit Steinen}
		\label{Steinebefüllen}
	\end{figure}
	\item Schließen Sie den Netzstecker an eine Steckdose an. Das System ist nun betriebsbereit.
	\begin{figure}[H]
		\centering
		\resizebox{0.75\textwidth}{!}{%
			\begin{circuitikz}
				\tikzstyle{every node}=[font=\LARGE]
				\draw [line width=1pt, short] (-12.5,10.5) -- (6.25,10.5);
				\draw [line width=1pt, short] (6.25,10.5) -- (18.75,3);
				\draw [line width=1pt, short] (6.25,10.5) -- (6.25,29.25);
				\draw [line width=1pt, short] (8.75,13) -- (8.75,17);
				\draw [line width=1pt, short] (8.75,17) -- (12,15);
				\draw [line width=1pt, short] (12,15) -- (12,11);
				\draw [line width=1pt, short] (8.75,13) -- (12,11);
				\draw [line width=1pt, short] (8.75,17) -- (9,17);
				\draw [line width=1pt, short] (12,15) -- (12.25,15);
				\draw [line width=1pt, short] (12,11) -- (12.25,11);
				\draw [line width=1pt, short] (12.25,15) -- (12.25,11);
				\draw [line width=1pt, short] (9,17) -- (12.25,15);
				\draw [ line width=1pt ] (10.5,14) ellipse (1cm and 1.5cm);
				\draw [line width=1pt, short] (11,12.75) .. controls (10.25,13.25) and (10.25,14.5) .. (11,15.25);
				\draw [line width=1pt, short] (11,14.25) .. controls (10.75,14.25) and (10.75,13.5) .. (11,13.75);
				\draw [line width=1pt, short] (11,14.25) .. controls (11.25,14) and (11.25,13.75) .. (11,13.75);
				\draw [line width=1pt, short] (-2,15.75) -- (0,15.75);
				\draw [line width=1pt, short] (-2,13.5) -- (0,13.5);
				\draw [line width=1pt, short] (0,15.75) -- (0,13.5);
				\draw [line width=1pt, short] (-2,16) -- (-0.25,16);
				\draw [line width=1pt, short] (-0.25,16) -- (0,15.75);
				\draw [line width=1pt, short] (0,15.5) -- (1.25,15.5);
				\draw [line width=1pt, short] (1.25,15.5) -- (1.25,15.25);
				\draw [line width=1pt, short] (1.25,15.25) -- (0,15.25);
				\draw [line width=1pt, short] (0,14) -- (1.25,14);
				\draw [line width=1pt, short] (1.25,14) -- (1.25,13.75);
				\draw [line width=1pt, short] (1.25,13.75) -- (0,13.75);
				\draw [line width=1pt, short] (1.25,15.5) -- (1,15.75);
				\draw [line width=1pt, short] (1,15.75) -- (0,15.75);
				\draw [line width=1pt, short] (1.25,14) -- (1,14.25);
				\draw [line width=1pt, short] (1,14.25) -- (0,14.25);
				\draw [line width=1pt, short] (-2,15.75) .. controls (-3.75,15.5) and (-3.5,13.5) .. (-2,13.5);
				\draw [line width=1pt, short] (-2,16) .. controls (-2.75,16) and (-3.25,15.25) .. (-3.25,14.75);
				\draw [line width=1pt, short] (2.5,15.5) -- (3.75,15.5);
				\draw [line width=1pt, short] (3.75,15.5) -- (3.75,16.25);
				\draw [line width=1pt, short] (3.75,16.25) -- (5,15);
				\draw [line width=1pt, short] (5,15) -- (3.75,14);
				\draw [line width=1pt, short] (3.75,14) -- (3.75,14.75);
				\draw [line width=1pt, short] (3.75,14.75) -- (2.5,14.75);
				\draw [line width=1pt, short] (2.5,14.75) -- (2.5,15.5);
				\draw [line width=1pt, short] (-3.25,15) .. controls (-9,15) and (-6.5,7.75) .. (-12.5,8);
				\draw [ line width=1pt ] (-12.5,3) rectangle (31.25,29.25);
			\end{circuitikz}
		}%
		\caption{Netzstecker anschließen}
		\label{Steckdose}
	\end{figure}
\end{enumerate}

